\section{Polars as a data worker: population, information time and execution meaning}
\label{sec:exec23-polars-worker}

A banking analysis can preserve every numerical formula and still change its
question when it joins customer records to events. A customer with three
contacts can acquire three times the weight of a customer with one contact;
an apparently historical measurement can have arrived after the decision;
and an overlapping time window can be mistaken for an independent
observation. A dataframe engine is valuable because it makes these
transformations explicit and executable. It cannot decide which
transformation represents the banking experiment.

Polars is therefore considered here as a conditional data-worker reference.
Twenty-one selected files authenticate to revision
\texttt{bf2152618c7123836a825f8566f4d98585b6ff5b}, dated 7 October 2026 before
the Asia/Shanghai cutoff. The Python project declares version 2.0.0,
Python at least 3.10 and the matching \texttt{polars-runtime-32} dependency;
the Rust workspace separately declares 0.55.1 and edition 2024. These are
source declarations, not authenticated distributed releases. Runtime-64,
compatibility, Arrow, pandas, database, cloud and GPU routes have distinct
dependencies. The selected root contains MIT terms, including notice
retention, but this is not clearance of every dependency or commercial
service. No downloaded implementation or test is executed in this review.
\citep{exec23_polars_sources}

\subsection{A join changes weights unless the population contract prevents it}

The selected Python lazy interface defaults to an inner join and
\texttt{validate="m:m"}; the latter performs no uniqueness check. It offers
one-to-one, one-to-many and many-to-one validation, and forwards validation,
null equality, coalescing and order choices to its underlying join call.
The guide distinguishes inner, outer, semi, anti, cross and as-of joins.
These are substantive alternatives: a semi join filters membership without
adding the right-hand columns, whereas an ordinary equi join can multiply
rows. They must be chosen from the observation design, not from the fact
that a table happens to fit in memory.\citep{exec23_polars_sources}

Let customer $i$ have outcome $Y_i$ and $m_i$ matching event records. For
one group with $n$ original customers and $\sum_i m_i>0$, the unweighted
mean after an ordinary join is
\begin{equation}
 \overline Y_{\mathrm{joined}}
 =\frac{\sum_{i=1}^n m_iY_i}{\sum_{i=1}^n m_i}
 =\overline Y+\frac{\operatorname{Cov}_n(m,Y)}{\overline m},
 \label{eq:exec23-polars-fanout}
\end{equation}
where $\operatorname{Cov}_n$ uses denominator $n$ and customers with no
match have $m_i=0$. The identity follows by expanding
$n^{-1}\sum_i(m_i-\overline m)(Y_i-\overline Y)$. Equal multiplicities
preserve the mean; dependence between outcome and multiplicity generally
does not. Neither the identity nor a successful join establishes a causal
effect.\citep{exec23_polars_reference}

For an exact supplied example, two treated customers have outcomes 10 and
0 and two controls have outcomes 0 and 0. Their customer-level descriptive
contrast is 5. If the first treated customer has three event records and
each other customer has one, the joined treated outcomes are
$(10,10,10,0)$ and the contrast becomes $15/2$. Against a stipulated
cost threshold of USD 6, the two descriptive calculations imply margins
of USD $-1$ and USD $3/2$. That reversal illustrates a population error;
it is not evidence that either contrast is a treatment effect or that
Polars has produced an incorrect join. An event-level target may be
appropriate, but it is a different target.\citep{exec23_polars_reference}

The proposed worker should carry a sampling-unit key, an expected join
cardinality and pre/post membership counts. A many-to-one join requires a
justified unique right-hand key. If multiple rows are meaningful revisions,
aggregation or version selection must be scientifically specified before
joining. Dropping duplicates with an arbitrary first-row rule merely
hides the unresolved choice. Anti joins can expose unmatched IDs for
review, while a declared refusal can preserve an unresolved key rather
than silently removing the customer.

\subsection{Missing values and missing identities are different problems}

Polars distinguishes missing \texttt{null} values from floating-point
\texttt{NaN}. The selected guide states that \texttt{null\_count} excludes
NaN, \texttt{fill\_null} does not fill it, and numerical means and sums
skip nulls while NaN typically propagates. These are useful computational
contracts. They do not say that omitted outcomes are missing at random or
that a nonfinite number should be reclassified as an unobserved bank
response.\citep{exec23_polars_sources}

Suppose the recorded outcomes are 2, 4 and one missing value. The observed
mean is 3. Filling the missing entry with zero changes the three-person
mean to 2. Two possible complete populations, with hidden values 0 and 8,
have the same observed data but full means 2 and $14/3$. Thus the
computational choice to skip, fill or reject does not identify the
full-population mean. A separate observation model or sensitivity range
must justify an estimand involving the missing response. These are exact
authored finite populations, not a claim about a package's imputation
accuracy.\citep{exec23_polars_reference}

Join keys require another distinction. The selected interface defaults to
not matching null keys, while \texttt{nulls\_equal=True} permits them to
match. An inspected unit test describes two null keys on the left matching
a null key on the right; it has not been run here. Equality of missing
tokens is a declared relation on values, not evidence that two unidentified
customers are the same person. The worker must preserve unmatched and
ambiguous identities and their denominators. Treating every missing
identifier as one customer would impose a linkage claim that the engine
has no information to establish.\citep{exec23_polars_sources}

\subsection{An as-of key must represent the information available to the decision}

The lazy as-of interface documents a backward search selecting the last
right-hand key at or before the left-hand key, a forward search selecting
the first key at or after it, and a nearest search. It defaults to
backward, permits exact matches unless told otherwise, and supports
grouping and a tolerance. Both inputs must be ordered within the relevant
groups. The same source qualifies checking: the in-memory engine cannot
check grouped sortedness, while the streaming engine checks the rows it
processes. Inspected tests describe these distinctions and deliberate
invalid sortedness declarations; they are intended checks, not observations
of this workspace's runtime.\citep{exec23_polars_sources}

Let $t$ be the decision time, $e_j$ a measurement's event time and $a_j$
its availability time. A possible historical feature contract chooses
the largest $e_j$ among records satisfying both $e_j\leq t$ and
$a_j\leq t$, within the correct customer and permitted revision lineage.
Backward matching on event time alone imposes only the first condition.
At a supplied decision time 100, consider records
$(e,a,x)=(90,105,0),(80,85,4),(101,101,10)$. Backward event matching
selects value 0, although it arrived too late; nearest matching selects
the future value 10. The stated availability-aware contract selects value
4. The exact example demonstrates information leakage from the wrong
clock, not a defect in nearest or backward matching themselves.
\citep{exec23_polars_reference}

Equal timestamps need a declared ordering convention: a value logged at
the decision timestamp may have been recorded just after the action.
Disabling exact matches can express strict inequality but cannot recover
an unrecorded within-timestamp order. Tolerances likewise need units and
meaning. The selected source distinguishes calendar-day tolerances from
fixed 24-hour durations under daylight saving. Its time-zone guide
distinguishes converting a zone from replacing one. A wrapper should bind
the zone, precision, strictness, revision and availability clock instead
of allowing a serialization round trip to choose them implicitly.
\citep{exec23_polars_sources}

Window construction also enters uncertainty. The guide permits the
window period to differ from its start interval, allowing overlap or
gaps. For independent mean-zero, variance-one signs $X_1,X_2,X_3$, define
$W_1=(X_1+X_2)/2$ and $W_2=(X_2+X_3)/2$. Then
$\operatorname{Var}(W_1)=\operatorname{Var}(W_2)=1/2$ and
$\operatorname{Cov}(W_1,W_2)=1/4$, so
\begin{equation}
 \operatorname{Var}\!\left(\frac{W_1+W_2}{2}\right)=\frac{3}{8}.
 \label{eq:exec23-polars-overlap}
\end{equation}
Treating the two windows as independent gives $1/4$, understating variance
by a factor of $3/2$. Enumeration of the eight sign triples reproduces
these moments. A correct window aggregation therefore still requires a
dependence-aware uncertainty calculation for the resulting panel or time
series.\citep{exec23_polars_sources,exec23_polars_reference}

\subsection{Streaming is a plan property, not a memory or recovery certificate}

The lazy interface exposes optimized and unoptimized plans. The guide
describes predicate, projection and slice pushdown, common-subplan
elimination, join ordering and type coercion. These can preserve the
declared relational query while reducing work; they do not determine
whether that query matches the experiment. The appropriate audit binds
the logical request and input identities as well as the selected physical
plan, engine and environment. An optimized graph is useful evidence of
execution choices, not evidence that a selected sample is causally valid.
\citep{exec23_polars_sources}

The selected \texttt{collect} documentation says that auto selection uses
configured engine affinity, otherwise streaming, and that unsupported
work can fall back to in-memory execution. GPU and background modes carry
explicit instability qualifications. The engine delegate forwards the
selected engine name to the underlying collect call; source comments
about relative performance are not measurements on this banking workload.
The streaming guide itself contains an editorial exclusion comment, so
its presence in the repository does not establish publication of that
page in the hosted manual.\citep{exec23_polars_sources}

A documentation tension is more informative when followed into code.
The eager join docstring says validation is not supported by the
streaming engine; the lazy docstring lists validation and streaming
support without that same note. In the inspected Rust lowering path,
the specialized streaming join branch requires
\texttt{!args.validation.needs\_checks()}. Other cases produce an
\texttt{InMemoryJoin} node while retaining the join arguments.
\texttt{needs\_checks} is false only for many-to-many validation.
The inspected in-memory node waits for both input sinks, invokes its
join closure and classifies its sink state as accumulating memory.
This supports a precise source-level interpretation: requesting checked
join cardinality can introduce an in-memory join within a streaming
plan. It does not support saying that validation is silently discarded,
that all such queries fail, or that native execution has been verified.
\citep{exec23_polars_sources}

This tradeoff matters to a proposal's cost model. Removing a key check to
retain streaming changes the accepted-input contract. Moving that check
to a separate prepass adds I/O and requires the checked snapshot to be
the one subsequently joined. The future trial must measure complete
validation, scans, conversion, spilling, fallback, output and restart
costs against an identical accepted-output definition. Neither a runtime
flag nor a successful small fixture supplies a bank memory bound or a
durable recovery guarantee.

\subsection{A narrow reuse choice with substantive alternatives}

The present choice is reference, with a Python process worker as a
candidate prototype route and direct Rust use as a separate engineering
alternative. Reusing a pinned engine can preserve its rich join and
expression semantics. Porting that engine into R, Java or C++ would add
compiler, null, ordering and dependency obligations without demonstrated
benefit. Consumers in those languages can instead exchange explicitly
typed artifacts through the common CLI or MCP contract; that bridge is
proposed and still needs conformance evidence.

A relational database or another dataframe engine remains a substantive
alternative when its snapshot, key and query contracts are better suited
to the chosen workload. No speed ranking is assigned without a matched
trial. Small project-owned enumerations remain the independent oracle
for population, clock and dependence examples. The scientific owner
specifies units, population, assignment and information law; the adapter
maintainer binds schema, null/NaN, keys, ordering and runtime versions;
the operator measures full resource and recovery behavior.

Before native integration, freeze actual Python and Rust runtime
artifacts, optional engines, dependency notices, input snapshots, plans
and output validation. Register edge cases for duplicate and missing
keys, nonfinite responses, late revisions, ties, grouped sortedness,
overlapping windows and checked-join fallback. Preserve failures and
refusals and compare full accepted-output costs. Source review supports
the reuse decision; native execution, causal qualification and bank
adoption remain open.
